\chapter{Install Jekyll}\label{chap: Install Jekyll}

As noted already, installing Jekyll on Linux is easy. Windows and Mac users need to follow the instructions at \url{https://jekyllrb.com/docs/}.

Linux Mint users, or others who have a Debian based system---Ubuntu for example---simply need to run the following command:

\begin{lstlisting}[style=myStyle,language=bash]
sudo apt install jekyll
\end{lstlisting}

You can also use your normal software installer application to install it if you prefer the GUI over the command line.

Once installed, execute the following commandd to get the version installed:

\begin{lstlisting}[style=myStyle,language=bash]
jekyll --version
\end{lstlisting}

Mine says ``4.3.2'' which is not the latest available, but is suitable for our use. Another optional test, to ensure that you have the same modules as I do, is to run this next command:

\begin{lstlisting}[style=myStyle,language=bash]
apt search jekyll | grep ^i
\end{lstlisting}

The output should list the following modules. Any that are not present will need installing individually with \code{sudo apt install <whatever>}:

\begin{lstlisting}[style=myText]
i   jekyll - simple, blog aware, static site generator
i A jekyll-theme-minima - beautiful, minimal theme for jekyll
i A ruby-jekyll-feed - Atom feed generator for Jekyll
i A ruby-jekyll-sass-converter - Basic Sass converter for Jekyll
i A ruby-jekyll-seo-tag - Jekyll plugin to add metadata tags
i A ruby-jekyll-watch - Rebuild your Jekyll site when a file changes
\end{lstlisting}

That's the installation done. Time to set up a small test website, only on your local PC, to get some practice.
